\documentclass[10pt, aspectratio=169]{beamer}

\usetheme{Madrid}
\usecolortheme{dolphin}

\usepackage{ctex}
\usepackage{xeCJK}
\usepackage{amsmath, amssymb}
\usepackage{graphicx}
\usepackage{listings}
\usepackage{tikz}
\usepackage{xcolor}

\title{408考研零基础 --- 数据结构}
\subtitle{1-12 线索二叉树}
\author{}
\date{}

\setbeamertemplate{page number in head/foot}{}

\begin{document}


\begin{frame}{本节目标}
    \begin{enumerate}
        \item 理解线索化的动机与原理
        \item 掌握二叉树的中序线索化算法
        \item 掌握线索二叉树中查找前驱与后继的方法
    \end{enumerate}
\end{frame}

\begin{frame}{线索化的动机}
    \textbf{$n$ 个结点的二叉链表}
    \begin{itemize}
        \item 共 $2n$ 个指针域
        \item $n-1$ 个指向孩子
        \item \textbf{$n+1$ 个空指针域}
    \end{itemize}
    \vspace{0.5em}
    \textbf{线索化思想}
    \begin{itemize}
        \item 空左指针 $\to$ 指向前驱
        \item 空右指针 $\to$ 指向后继
        \item 增加标志位区分孩子/线索
    \end{itemize}
\end{frame}

\begin{frame}{线索二叉树的存储结构}
    \begin{lstlisting}[language=C, basicstyle=\ttfamily\small]
typedef struct ThreadNode {
    ElemType data;
    struct ThreadNode *lchild, *rchild;
    int ltag, rtag;
} ThreadNode, *ThreadTree;
    \end{lstlisting}
    \vspace{0.5em}
    \textbf{标志位含义}
    \begin{itemize}
        \item $\text{ltag} = 0$：lchild 指向左孩子
        \item $\text{ltag} = 1$：lchild 指向前驱
        \item $\text{rtag} = 0$：rchild 指向右孩子
        \item $\text{rtag} = 1$：rchild 指向后继
    \end{itemize}
\end{frame}

\begin{frame}{中序线索化算法}
    \textbf{基于中序遍历，使用 pre 指针}
    \vspace{0.3em}
    \begin{enumerate}
        \item 线索化左子树
        \item 处理当前结点 $p$
        \begin{itemize}
            \item 若 $p \to \text{lchild} == \text{NULL}$，指向 pre
            \item 若 $\text{pre} \to \text{rchild} == \text{NULL}$，指向 $p$
        \end{itemize}
        \item $\text{pre} = p$
        \item 线索化右子树
    \end{enumerate}
    \vspace{0.3em}
    时间复杂度 $O(n)$，空间复杂度 $O(1)$
\end{frame}

\begin{frame}{中序线索二叉树的遍历}
    \textbf{找中序后继}
    \begin{itemize}
        \item $\text{rtag} = 1$：直接 $p \to \text{rchild}$
        \item $\text{rtag} = 0$：右子树中最左边的结点
    \end{itemize}
    \vspace{0.3em}
    \textbf{找中序前驱}
    \begin{itemize}
        \item $\text{ltag} = 1$：直接 $p \to \text{lchild}$
        \item $\text{ltag} = 0$：左子树中最右边的结点
    \end{itemize}
    \vspace{0.3em}
    \textbf{遍历算法}
    \begin{enumerate}
        \item 找到中序第一个结点（沿左走到 ltag=1）
        \item 循环调用找后继，访问每个结点
    \end{enumerate}
    时间复杂度 $O(n)$，空间复杂度 $O(1)$
\end{frame}

\begin{frame}{本节总结}
    \begin{enumerate}
        \item 线索化利用 $n+1$ 个空指针记录前驱后继
        \item 中序线索化基于中序遍历，使用 pre 指针
        \item 线索二叉树遍历无需递归或栈
        \item 时间 $O(n)$，空间 $O(1)$
    \end{enumerate}
\end{frame}

\end{document}
